\bookchapter{פרק 0: היי, לפני שמתחילים}

לפני שנתחיל לחקור באופן רשמי, בואו נדבר על כמה כללים בסיסיים, כדי לוודא שתבין כל מה שייאמר בהמשך. אל תמהר, השחזת הסכין לא מעכבת את כריתת העצים.

\booksection{א. מה הספר הזה יכול לעזור לך? (ומה לא)}

הוא יכול לעזור לך להבין באופן שיטתי איך להשתמש בטלפון, להבין את ההיגיון מאחורי הפונקציות, ובין אם אתה איש אנדרואיד או אפל, רוב הרעיונות אוניברסליים. הוא ייקח אותך לתרגול מעשי, ויעלה אותך מ“יודע להשתמש” ל“מבין איך להשתמש”.

הוא אינו ספר ההוראות של טלפון מסוים, אינו חדשות טכנולוגיה על מוצרים חדשים, ובוודאי שאינו מדריך להתמכרות לטלפון. להיות כנים, אם ציפית ל“איך לגלול סרטונים בכיף יותר”, אולי תתאכזב - אנחנו יותר אכפת לנו איך לגרום לטלפון לשרת אותך באמת, ולא להיפך.

\booksection{ב. הסבר על “הצופן” שלנו}

כדי להסביר בבירור רב יותר, הספר משתמש בכמה סימנים קטנים, שהם כמו צופן מוסכם בינינו:

\begin{itemize}
  \item \textbf{[שם כפתור]}: מייצג כפתור ספציפי על המסך או שם App, למשל [הגדרות], [אישור].
  \item \textbf{>}: מייצג את סדר הפעולות, למשל "לחץ [הגדרות] > [צליל]", כלומר קודם לחץ על הגדרות, ואז על צליל.
  \item \textbf{מודגש}: חשוב מאוד! עשוי להיות מושג מפתח או אזהרת בטיחות, נא לשים לב.
  \item \emph{נטוי}: מחשבות משלימות קטנות או מטאפורות, לעזור לך להבין טוב יותר.
\end{itemize}

\booksection{ג. אנדרואיד ו-iOS, נדבר על שניהם}

תנחש מה, יש דברים שנקראים בשם שונה אך עושים אותו דבר. למשל: יישום (אנדרואיד) ו-App (iOS), הגדרות (שניהם נקראים הגדרות), ניהול קבצים (אנדרואיד) וקבצים (iOS). העיקר שתדע שזה אותו דבר, אל תתבלבל מהמונחים.

בנוסף, כיום במחנה אנדרואיד יש גם את מערכת ההרמוני של Huawei, שיש לה את המאפיינים שלה, נדבר עליה במיוחד בהמשך. בקיצור, בין אם אתה משתמש באיזה טלפון שהוא, הספר הזה יכול לעזור לך.

\booksection{ד. קיבולת הטלפון: למה קנית 128G אך השטח הפנוי פחות מ-110G?}

אנשים מכירים את יחידות הקיבולת: 1TB = 1024GB, 1GB = 1024MB, 1MB = 1024KB... אבל יש כאן מלכודת: יצרני דיסקים בדרך כלל מחשבים בבסיס 1000, בעוד המערכת מחשבת בבסיס 1024, לכן הדיסק “128G” שקנית בפועל “יתכווץ”.

אך הסיבות העיקריות לכיווץ הקיבולת של הטלפון שלך הן שלוש:

\begin{enumerate}
  \item \textbf{המערכת עצמה תופסת מקום}: מערכת ההפעלה היא כמו קירות נושאים של בית, חייבת להתקיים. המערכות של ימינו תופסות לעיתים קרובות 10-20GB, זה נורמלי.
  \item \textbf{יישומים מותקנים מראש על ידי היצרן}: ה-App שמגיעים עם הטלפון החדש גם תופסים מקום, חלקם ניתנים להסרה וחלקם לא.
  \item \textbf{חלוקה ופורמט איבוד}: סיבות טכניות, כמו שבד הריפוד של תיק בית הספר גם תופס נפח.
\end{enumerate}

\begin{tcbinfo}{טיפ קטן לקניית טלפון}
כיום טלפון 128GB כבר לא מספיק בעצם, 256GB ומעלה הוא התחלה יציבה יותר. אם אתה אוהב לצלם ולשחק, מומלץ לעבור ישירות ל-512GB או אפילו יותר. זכור, שטח אחסון, עדיף עודף מאשר מחסור.

\end{tcbinfo}

\begin{tcbpractice}{בוא נעשה: ראה את "תעודת הזהות" של הטלפון שלך}
פתח את [הגדרות], מצא אודות הטלפון (או אפשרות דומה), רשום: דגם הטלפון, גרסת מערכת, שטח אחסון כולל ושטח פנוי. זהו הצעד הראשון להכיר אותו. בדוק גם את בריאות הסוללה, כדי שתהיה לך תמונה בראש.

\end{tcbpractice}

\booksection{ה. הטלפון שלך הוא “מאגר נכסים דיגיטליים”}

לפני שנתחיל בידע ספציפי, בואו נבנה הכרה מרכזית:\textbf{הטלפון כבר מזמן אינו רק כלי תקשורת, הוא מאגר הנכסים הדיגיטליים החשוב ביותר שלנו}. אנשי קשר, היסטוריית צ'אט, תמונות משפחה, מסמכי עבודה, App בנקאות, קוד תשלום, נתוני בריאות, תעודות אלקטרוניות... הנתונים האלה יחד, הערך שלהם לעיתים קרובות עולה על הטלפון עצמו.

\begin{tcbinfo}{למה זה חשוב?}
אם הטלפון אבד או נהרס, ההפסד שלך אינו רק אלפי שקלים של חומרה,\textbf{אלא יותר כנראה נתונים בלתי מתחדשים שנצברו במשך שנים}- תמונות הגדילה של הילדים, הקלטות קול של קרובים שנפטרו, חוזים חשובים שלא גובו. הרגלי ארגון והגנה מדעיים יכולים להגן עליך בשלושה ממדים:

\begin{itemize}
  \item \textbf{שיפור יעילות פעולה}: סיווג קבצים תקני מאפשר למערכת למצוא מהר יותר, הטלפון לא נתקע לאורך זמן.
  \item \textbf{הארכת חיי שימוש}: ניהול אחסון סביר מפחית את עומס הקריאה והכתיבה של השבב, הסוללה והאחסון עמידים יותר.
  \item \textbf{הבטחת אבטחת נתונים}: ניהול שכבתי + גיבוי תקופתי + צמצום הרשאות, מרחיק "אסון נתונים" ממך.
\end{itemize}

\end{tcbinfo}

\booksubsection{הרעיון המרכזי שלנו: “מניעה בראש, שילוב מניעה וטיפול”}

כמו שלא מחכים לחלות כדי ללכת לרופא,\textbf{ניהול הטלפון צריך להתחיל מ“הרגלים יומיומיים”}. הספר הזה מדגיש שוב ושוב לא טכניקה מתקדמת כלשהי, אלא את העקרונות הפשוטים האלה:

\begin{itemize}
  \item \textbf{מניעה בראש}: התקנת App רק דרך חנות רשמית, הסרת רגישות לפני שליחת תמונות, לא לחזור על סיסמאות - ה"פעולות הקטנות" האלה בעלות נמוכה מאוד, אך יכולות לחסום את רוב הסיכונים.
  \item \textbf{בדיקה תקופתית}: סידור תיקיית ההורדות פעם בחודש, ביקורת הרשאות פעם ברבעון, גיבוי מלא פעם בחצי שנה - קצב קבוע יעיל פי מאה מכיבוי שריפות אד-הוק.
  \item \textbf{להשאיר דרך חזרה}: הפעלת סנכרון ענן, רישום חשבונות וסיסמאות, כתיבת "רשימת מורשת דיגיטלית" - ודא שבעת תקלה גם אתה וגם משפחתך תוכלו להתמודד בשלווה.
\end{itemize}
